\section{Term calculus}
\label{sec:calculus}

The source language is the Calculus of Constructions extended with inductive types and fixpoints. Terms follow the grammar:
\[
\begin{array}{rcl}
 t, A, B &::=& \Var{x} \mid \Sort{i} \mid \PiTy{A}{B} \mid \Lam{A}{t}
                \mid \App{t}{u} \mid \Fix{A}{t} \\
         &\mid& \Ind{I}{\vec{a}} \mid \Roll{I}{c}{\vec{a}}
                \mid \Case{I}{t}{C}{\vec{b}}
\end{array}
\]
where $x$ ranges over de~Bruijn indices, $\Sort{i}$ is the universe hierarchy,
$\PiTy{A}{B}$ is the dependent function type, $\Lam{A}{t}$ is lambda
abstraction, $\App{t}{u}$ is application, $\Fix{A}{t}$ is a recursive
definition (recursive call bound at de~Bruijn index 0), $\Ind{I}{\vec{a}}$
is an applied inductive type, $\Roll{I}{c}{\vec{a}}$ is a constructor
application, and $\Case{I}{t}{C}{\vec{b}}$ is dependent pattern matching
with motive $C$.  Contexts $\ctx$ are lists of types; the global
environment $\env$ records inductive definitions.

\paragraph{Typing (compact).}
The judgement $\hastype{\env}{\ctx}{t}{A}$ is standard for CIC with
inductives~\cite{coquand-huet-1988-coc,paulin-mohring-1993-inductive}.
Conversion uses $\beta$-convertibility.  The core rules (Figure~\ref{fig:typing})
are:

\bigskip
\noindent
\begin{minipage}{0.48\linewidth}
\begin{prooftree}
  \AxiomC{$\ctx(x)=A$}
  \RightLabel{\scriptsize (TyVar)}
  \UnaryInfC{$\hastype{\env}{\ctx}{\Var{x}}{A}$}
\end{prooftree}
\end{minipage}%
\begin{minipage}{0.48\linewidth}
\begin{prooftree}
  \AxiomC{}
  \RightLabel{\scriptsize (TySort)}
  \UnaryInfC{$\hastype{\env}{\ctx}{\Sort{i}}{\Sort{i+1}}$}
\end{prooftree}
\end{minipage}

\bigskip
\begin{minipage}{0.48\linewidth}
\begin{prooftree}
  \AxiomC{$\hastype{\env}{\ctx}{A}{\Sort{i}}$}
  \AxiomC{$\hastype{\env}{A{::}\ctx}{B}{\Sort{j}}$}
  \RightLabel{\scriptsize (TyPi)}
  \BinaryInfC{$\hastype{\env}{\ctx}{\PiTy{A}{B}}{\Sort{\max(i,j)}}$}
\end{prooftree}
\end{minipage}%
\begin{minipage}{0.48\linewidth}
\begin{prooftree}
  \AxiomC{$\hastype{\env}{\ctx}{A}{\Sort{i}}$}
  \AxiomC{$\hastype{\env}{A{::}\ctx}{t}{B}$}
  \RightLabel{\scriptsize (TyLam)}
  \BinaryInfC{$\hastype{\env}{\ctx}{\Lam{A}{t}}{\PiTy{A}{B}}$}
\end{prooftree}
\end{minipage}

\bigskip
\begin{minipage}{0.48\linewidth}
\begin{prooftree}
  \AxiomC{$\hastype{\env}{\ctx}{t}{\PiTy{A}{B}}$}
  \AxiomC{$\hastype{\env}{\ctx}{u}{A}$}
  \RightLabel{\scriptsize (TyApp)}
  \BinaryInfC{$\hastype{\env}{\ctx}{\App{t}{u}}{B[u/]}$}
\end{prooftree}
\end{minipage}%
\begin{minipage}{0.48\linewidth}
\begin{prooftree}
  \AxiomC{$\hastype{\env}{\ctx}{A}{\Sort{i}}$}
  \AxiomC{$\hastype{\env}{A{::}\ctx}{t}{A\uparrow}$}
  \RightLabel{\scriptsize (TyFix)}
  \BinaryInfC{$\hastype{\env}{\ctx}{\Fix{A}{t}}{A}$}
\end{prooftree}
\end{minipage}

\bigskip
\begin{prooftree}
  \AxiomC{$\hastype{\env}{\ctx}{t}{\Ind{I}{\vec{\delta}}}$}
  \AxiomC{$\hastype{\env}{\ctx,\,x{:}\Ind{I}{\vec{\delta}}}{C}{\Sort{i}}$}
  \AxiomC{$\hastype{\env}{\ctx}{b_c}{\Pi\vec{\tau}_c.\,C[\Roll{I}{c}{\vec{y}}/x]}$ per $c$}
  \RightLabel{\scriptsize (TyCase)}
  \TrinaryInfC{$\hastype{\env}{\ctx}{\Case{I}{t}{C}{\vec{b}}}{C[t/x]}$}
\end{prooftree}

\noindent plus (TyInd) and (TyRoll) which follow the inductive family
structure~\cite{paulin-mohring-1993-inductive} as formalised in the
mechanisation.

\paragraph{Operational semantics.}
We use call-by-name (CBN) weak-head reduction with three rules:
\[
\App{(\Lam{A}{t})}{u} \step t[u/] \qquad
\Fix{A}{t} \step t[\Fix{A}{t}/] \qquad
\Case{I}{(\Roll{I}{c}{\vec{a}})}{C}{\vec{b}} \step b_c\,\vec{a}
\]
plus congruence for reduction in the function position of application and
the scrutinee of case expressions.  Under CBN, arguments are never
evaluated before substitution; reduction does not proceed under binders.
The reflexive-transitive closure is written $\steps$.

\paragraph{Observational equivalence.}
We use CIU (Closed Instantiations of Uses) equivalence: $t \approx u$ iff
for all closing substitutions $\sigma$ and CBN values $v$, $t[\sigma]$
terminates to $v$ exactly when $u[\sigma]$ does.  This captures the standard
notion of contextual equivalence for CBN and is the correctness criterion
for our supercompilation pipeline (Claim~3).
